\addtocontents{toc}{\protect\Needspace{3\baselineskip}}
\originalchapter{场的振子与粒子态}\label{ch:basic-particles}
对应 8.251 第 16 讲后半、第 17 讲. 经典场中的平面波不是自动等于一个粒子; 正则量子化后, 创建算符作用于真空产生能量与动量本征态, 才赋予粒子解释. 这里只补足弦谱所需的自由场结构, 更广泛的相互作用与重整化以前置 QFT 卷为准.

\originalsection{实标量场的方程与能量}
自然单位下质量记为 $m$, 作用量为
\[
 S=\int\dd^Dx\left[-\frac12\partial_\mu\phi\partial^\mu\phi
 -\frac12m^2\phi^2\right].
\]
变分并令外边界贡献为零, 得 $(\partial^2-m^2)\phi=0$, 即
$(-\partial_t^2+\nabla^2-m^2)\phi=0$. 原课笔记个别展开式把质量项符号写成正号; 本稿用上述作用量与方程保持一致. 共轭动量 $\pi=\dot\phi$, 能量密度
\[
 \mathcal H=\frac12\pi^2+\frac12(\nabla\phi)^2+\frac12m^2\phi^2.
\]
若 $m^2\geq0$, 经典能量非负; 若 $m^2<0$, 是关于所展开背景的稳定性问题, 后面讨论弦的快子时会用到.

\originalsection{平面波与实条件}
试解 $\phi\propto\ee^{-\ii Et+\ii\boldsymbol p\cdot\boldsymbol x}$ 得
$E^2=\boldsymbol p^2+m^2$. 实场不能只保留一个复平面波, 要加共轭项. 完整傅里叶分布满足
\[
 \phi(-p)=\phi(p)^*,\qquad(p^2+m^2)\phi(p)=0.
\]
后一式的非零分布支持在质量壳上, 正负能量两支由现实性联系. 不能把 $\phi(p)$ 当任意普通连续函数并说它仅在零体积曲面上非零而不解释分布. 有效的写法是直接对空间动量积分, 将能量固定为 $E_{\boldsymbol p}=\sqrt{\boldsymbol p^2+m^2}$.

\originalsection{有限体积内的独立振子}
先把空间放在体积 $V=\prod_iL_i$ 的周期盒, $p_i=2\pi n_i/L_i$. 对每一对非零 $\boldsymbol p,-\boldsymbol p$, 实场的两个实 Fourier 系数可写成余弦、正弦两个坐标, 代入作用量得到频率均为 $E_{\boldsymbol p}$ 的谐振子. 这避免把同一实场的共轭系数重复当作两份独立复自由度.

量子场展开统一写为
\[
 \phi(t,\boldsymbol x)=\frac1{\sqrt V}\sum_{\boldsymbol p}
 \frac{a_{\boldsymbol p}\ee^{-\ii E_{\boldsymbol p}t+\ii\boldsymbol p\cdot\boldsymbol x}
 +a_{\boldsymbol p}^\dagger\ee^{\ii E_{\boldsymbol p}t-\ii\boldsymbol p\cdot\boldsymbol x}}
 {\sqrt{2E_{\boldsymbol p}}},
\]
其中 $[a_{\boldsymbol p},a_{\boldsymbol q}^\dagger]=\delta_{\boldsymbol p\boldsymbol q}$, 其余对易为零. 求 $\pi=\dot\phi$ 并代入, 一项与其共轭各给半份核,
\[
 [\phi(t,\boldsymbol x),\pi(t,\boldsymbol y)]
 =\frac\ii V\sum_{\boldsymbol p}\ee^{\ii\boldsymbol p\cdot(\boldsymbol x-\boldsymbol y)}
 =\ii\delta_V(\boldsymbol x-\boldsymbol y).
\]
因而 $1/\sqrt{2E}$ 归一化并非随意.

质量为零的空间零动量模式有 $E=0$, 不适用这个除 $\sqrt E$ 的振子展开, 应单独按自由坐标处理. 对连续无限体积的一粒子散射态, 这一个特殊零模式不改变通常质量壳结构.

\originalsection{创建算符为什么产生粒子}
将展开代入能量与动量, 并对无相互作用的真空常数正规序, 得
\[
 H=\sum_{\boldsymbol p}E_{\boldsymbol p}a_{\boldsymbol p}^\dagger a_{\boldsymbol p},
 \qquad\boldsymbol P=\sum_{\boldsymbol p}\boldsymbol p\,
 a_{\boldsymbol p}^\dagger a_{\boldsymbol p}.
\]
定义 $a_p|\Omega\rangle=0$. 使用 $[H,a_p^\dagger]=E_pa_p^\dagger$,
$[\boldsymbol P,a_p^\dagger]=\boldsymbol p a_p^\dagger$, 得
$a_p^\dagger|\Omega\rangle$ 有能量 $E_p$、动量 $p$、质量 $m$. 多个创建算符给相应能量动量的和. 它们可交换, 多粒子波函数具有玻色交换对称性.

自由场的真空能常数可以在非引力能差问题中选基准减去. 弦质量的截距却不能任意照搬这一做法, 因纵向哈密顿量与洛伦兹生成元的约束会限制它. 这是下一章需要单独检查的一致性条件.

\originalsection{光锥粒子与规范补偿}
对正能量粒子, $(p^+,p^I)$ 唯一决定 $p^-=(p^Ip^I+m^2)/(2p^+)$. 用 $x^+$ 演化时, 约化变量可取 $x_0^-,p^+,x_0^I,p^I$, 并保留 $[x_0^-,p^+]=-\ii$.
选择规范 $x^+(\tau)=\tau$ 时, 自由世界线为
\[
 x^I(\tau)=x_0^I+\frac{p^I}{p^+}\tau,\qquad
 x^-(\tau)=x_0^-+\frac{p^-}{p^+}\tau.
\]
这套粒子参数不是前章开弦的 $X^+=2\alpha'p^+\tau$. 洛伦兹变换可能改变 $x^+$, 所以必须补上世界线重参数化才能回到 $x^+=\tau$. 总的规范固定变换因此既有通常矢量变换, 也有沿世界线的移动. 习题 HW9.5 将用这些明确的约定核对其对易式, 而不是对一个未声明时间参数的轨迹猜补偿项.
